\section{Metodología}
El presente proyecto se desarrolló bajo un enfoque cualitativo, orientado a la comprensión e interpretación de conceptos vinculados con Git y GitHub, la Guía de Scrum 2020, los patrones de diseño, el patrón Modelo–Vista–Controlador (MVC) y el lenguaje Java. Desde el inicio de la investigación se definió que estas tecnologías y metodologías representarían la base teórica para proponer buenas prácticas en el desarrollo de software colaborativo y escalable. Con el fin de sustentar y validar estas elecciones, se llevó a cabo un proceso de búsqueda, selección y análisis de literatura académica y técnica que ofreciera evidencia de sus beneficios y casos de uso. Se consideró fundamental identificar no solo la aplicabilidad individual de cada tema, sino también las sinergias que surgen al integrarlos en proyectos de diversa magnitud.

\subsection{Fase 1 - Recolección y construcción del thesaurus}
• Se utilizó como fuente principal Google Scholar, complementada con repositorios de documentación oficial y publicaciones especializadas. Se realizaron consultas con palabras clave relacionadas: Git, GitHub, Scrum 2020, patrones de diseño, MVC y Java.

• Cada recurso relevante fue documentado en un archivo "Thesaurus" con campos estandarizados que incluían término, fecha, fuente y nombre del artículo o guía consultada.

• El objetivo de esta fase fue reunir evidencia científica y técnica que respaldara el uso de estas herramientas y enfoques, proporcionando una base sólida para las recomendaciones metodológicas y prácticas presentadas en este trabajo.

\subsection{Fase 2 - Análisis cualitativo de la evidencia (síntesis del thesaurus)}
• El objetivo de esta fase fue transformar la información recopilada en argumentos sólidos que justificaran la elección de Git y GitHub, la Guía de Scrum 2020, los patrones de diseño, el patrón MVC y el lenguaje Java como elementos clave en el desarrollo de software. Para ello se realizaron lecturas críticas, resúmenes comparativos y reflexiones sobre cada fuente, evaluando aspectos como la colaboración distribuida, la trazabilidad de cambios, la modularidad de las arquitecturas y la portabilidad del código.

• La evidencia analizada confirmó la pertinencia de integrar estas tecnologías y metodologías en un mismo proyecto, destacando que su uso conjunto favorece la escalabilidad, la calidad del producto final y la eficiencia del trabajo en equipo. Asimismo, proporcionó lineamientos prácticos que orientan tanto el diseño de la arquitectura como la implementación y el mantenimiento de sistemas modernos.